% Generated draft: requires mathematical and source review.
\documentclass[11pt]{article}
\usepackage[margin=25mm]{geometry}
\usepackage{fontspec}
\setmainfont{Latin Modern Roman}
\usepackage{amsmath,amssymb,mathtools,unicode-math}
\setmathfont{Latin Modern Math}
\usepackage{enumitem,needspace,xurl,hyperref}
\setlength{\parindent}{0pt}
\setlength{\parskip}{5pt}
\setlength{\emergencystretch}{3em}
\newcommand{\N}{\mathbb{N}}
\newcommand{\Z}{\mathbb{Z}}
\newcommand{\Q}{\mathbb{Q}}
\newcommand{\R}{\mathbb{R}}
\newcommand{\C}{\mathbb{C}}
\newcommand{\F}{\mathbb{F}}
\newcommand{\A}{\mathbb{A}}
\newcommand{\PP}{\mathbb{P}}
\newcommand{\E}{\mathbb{E}}
\newcommand{\eps}{\varepsilon}
\newcommand{\dd}{\,\mathrm{d}}
\begin{document}
\section*{UnsolvedMath ID: 30000409}
\textbf{Source problem code:} catalogue entry
\subsection{Definitions and mathematical statement}
\textbf{Global iterative algebraicity question.} Is every global iterative differential module \((M,\partial_M)\) over a global iterative differential ring \((\mathcal O_{F'},\partial_F)\) algebraic? The words \emph{global} and \emph{algebraic} are essential source-specific qualifications. The record does not define \(F,F'\), the ring \(\mathcal O_{F'}\), its global arithmetic conditions, or the allowed trivializing extensions; it must not be read as a claim for arbitrary differential modules.\par For orientation, an iterative higher derivation on a commutative ring \(R\) is a sequence of additive maps \(\partial^{(n)}:R\to R\), \(n\ge0\), with \(\partial^{(0)}=\operatorname{id}\),\[\partial^{(n)}(ab)=\sum_{i+j=n}\partial^{(i)}(a)\partial^{(j)}(b),\qquad \partial^{(i)}\partial^{(j)}=\binom{i+j}{i}\partial^{(i+j)}.\] An iterative module structure has corresponding additive maps on the intended \(R\)-module \(M\), satisfying \(\partial_M^{(0)}=\operatorname{id}\), the same iterativity identity, and \(\partial_M^{(n)}(am)=\sum_{i+j=n}\partial^{(i)}(a)\partial_M^{(j)}(m)\). The usual finite-free or finite-projective rank hypotheses and the global extension of these maps must be recovered from the source; they are not fixed merely by writing these identities. In the displayed question the ring is named \(\mathcal O_{F'}\) but the derivation is indexed by \(F\), whose compatibility is also not specified.\par In a common differential-module sense, algebraicity asks for trivialization after a finite algebraic extension of the relevant fraction field: the extended module has a basis of horizontal vectors \(e\) with \(\partial_M^{(n)}e=0\) for every \(n\ge1\). Whether this is the source's definition, which constants must be preserved, whether separability is required, and whether trivialization is over a field or an integral ring are missing. One must not replace this with the much weaker statement that a differential Galois group is an algebraic group; the intended finite-extension or finite-group condition is different. No characteristic-zero factorial reduction of all higher derivations to one ordinary derivation is used without a stated characteristic hypothesis.\par The supplied triage checked 2026-08-21 labels the question open and identifies it as a global iterative form of Grothendieck's \(p\)-curvature conjecture. It says the report presents the statement as conjectural and no general proof was found in its search. The global hypotheses and the precise equivalence with a \(p\)-curvature formulation are not reproduced, so neither is asserted as a verified theorem here. This dated assessment requires independent verification.
\subsection{Short English statement}
Determine whether the source's global iterative differential modules become trivial after an appropriate finite algebraic extension. The arithmetic meaning of global, the fields and ring, and the definition of algebraicity are missing. Dropping those qualifiers would change the conjecture.
\subsection{Sources}
\begin{itemize}\item Local UnsolvedMath registry, permanent ID 30000409, code catalogue entry; \emph{Algebraicity of Global Iterative Differential Modules}, Oberwolfach Reports Open Problems.\item \emph{The Arithmetic of Fields (2006)}, Oberwolfach Reports, DOI 10.4171/owr/2006/06, pp. 317--378. \url{https://doi.org/10.4171/owr/2006/06}. Proposed year 2006; no proposer named.\item Supplied literature label \emph{original OWR report}: \url{https://doi.org/10.4171/owr/2006/06}.\item Supplied literature label \emph{primary report text}: \url{https://ems.press/content/serial-article-files/46036}.\item Literature triage checked 2026-08-21, status \texttt{open}, with the Grothendieck \(p\)-curvature attribution; record updated 2026-09-16. The \texttt{corrected\_verified} field concerns delimiter normalization only. Verify dated claims independently.\end{itemize}
\end{document}
